\appendix
\begin{table*}[t]
  \caption{Detailed RQ3 measurements. Under CompCert, Time and Kernel are the
  wrapper and unchecked kernel times in nanoseconds per call; Ratio is
  Time/Kernel. Under Clang, QD is the QD time and Ratio is Time/QD. Entries
  are medians [interquartile ranges] over paired trials. A dash denotes the
  single public wrapper for an error-free transformation.}
  \label{tab:rq3-performance-detail}
  \centering
  \setlength{\tabcolsep}{4pt}
  \begin{tabular}{@{}llrrrrrr@{}}
    \toprule
    & & \multicolumn{3}{c}{CompCert}
    & \multicolumn{3}{c}{Clang} \\
    \cmidrule(lr){3-5}\cmidrule(l){6-8}
    Operation & Wrapper
      & Time & Kernel & Ratio
      & Time & QD & Ratio \\
    \midrule
    \input{generated/rq3_combined_rows}\\
    \bottomrule
  \end{tabular}
\end{table*}

\section{Evaluation Details}
\label{app:evaluation-details}

We record supplementary details of the experimental setup in this appendix:
the construction of the generated inputs, the numerical reference
calculations, and the timing protocol. We report and interpret the
resulting measurements in \Cref{sec:evaluation}.

\subsection{Input Construction}

We generate inputs without consulting the runtime checks. At every base-two
exponent, we include positive and negative values equally often.

For 95\% of the double-word inputs, we use MPFR to construct a source value
with 106 significant bits and a nonzero final bit. We obtain the high
component by rounding this value to binary64 and the low component by rounding
the exact residual.
If the leading bit is at $2^e$, the final source bit is at $2^{e-105}$. At the
lower endpoint reported in \Cref{sec:evaluation-setup}, this final bit is
$2^{-1074}$; at the upper endpoint, rounding the high component remains finite.
For the remaining 5\%, a generated binary64 value becomes the high component
and the low component is zero. These endpoints belong to the test
construction and are not restrictions on \FenceDD{}.

The RQ2 summaries group all failure statuses for each wrapper. The artifact
records the exact status returned by every call, together with the run
identifier, random seed, input fingerprint, and component bit patterns.

\subsection{Numerical Reference Calculations}

We use QD 2.3.24, MPFR 4.2.2, and GMP 6.3.0. We enable QD's FMA
implementation of \ccode|two_prod|.

For RQ1, we use MPFR to reconstruct the real values represented by the input
and output pairs and to compute each reference operation at 2112-bit
precision. We recompute division cases whose reported error ratio or region
classification could depend on this precision at 4224-bit precision. We
confirm ratios close to one using directed rounding or interval calculations.
No result produced by an accepted RQ1 call has $B(r)=0$.

For each double-word operation, we select the QD routine computing the same
mathematical operation. For addition of two double words, we use
\ccode|ieee_add|; for division of two double words, we use
\ccode|accurate_div|; and for the remaining operations, we use the
corresponding QD operators. For the error-free transformations, we compare
with QD's
\ccode|two_sum|, \ccode|quick_two_sum|, and \ccode|two_prod|.

For RQ2, we classify results obtained using QD as nonfinite; finite with
components that do not form a double word; or finite with components that form
a double word, subdivided by whether the error exceeds the
corresponding \FenceDD{} bound. For an error-free transformation, the relevant
numerical property instead requires exact reconstruction: the returned
components must sum to the exact result.

\subsection{Timing Protocol}
\label{app:evaluation-timing}

Every wrapper returns successfully on the timing inputs, and before timing we
confirm that each wrapper returns the same result as its unchecked kernel. We
construct QD objects before measurement.

For the comparison with the kernels, we compile the \FenceDD{} wrappers and
kernels with CompCert 3.17 using \texttt{-O3}. For the comparison with QD, we
compile the \FenceDD{} C objects with Homebrew Clang 19.1.7 and QD and the C++
driver with the matching Homebrew Clang++ 19.1.7. For the Clang build, we use
\texttt{-O3}, with fast mathematics and implicit floating-point contraction
disabled. The test machine is a 16-core Apple M4 Max running macOS 26.6.2.

The 24 trials for each double-word operation use every ordering of the kernel,
default wrapper, alternative wrapper, and QD once. For each error-free
transformation, the six orderings of the kernel, wrapper, and QD each occur
four times. We give every implementation an untimed warm-up and compute
differences and ratios from paired trials.

\subsection{Detailed Performance Results}

We report the individual timing measurements and results for the alternative
wrappers in \Cref{tab:rq3-performance-detail}, supplementing the summary in
\Cref{tab:rq3-performance}.

Under CompCert, the alternative wrapper/kernel ratios range from
\RQThreeStrongMinimumKernelRatio{} to \RQThreeStrongMaximumKernelRatio{}, with
added times of \RQThreeStrongMinimumAddedTime{}--
\RQThreeStrongMaximumAddedTime{} nanoseconds per call. The Clang measurements
in \Cref{tab:rq3-performance-detail} compare the total running times of these
wrappers with QD.

\section{Proof of the Small-\FastTwoSum{} Result}
\label{app:fast-two-sum-small-proof}

\begin{proof}[Proof of \Cref{thm:fast-two-sum-small}]
Every value in the FLT format is an integer multiple of $\eta$, so the exact
sum $a+b$ lies on the same fixed lattice. The bound $|a+b|\leq2\nu$ ensures
that this lattice point is representable in the format. Line~1 of
\Cref{alg:fast-two-sum} therefore gives $s=a+b$. Consequently, Lines~2 and~3
give $z=\operatorname{RN}(s-a)=\operatorname{RN}(b)=b$ and
$t=\operatorname{RN}(b-z)=0$. Thus, $s+t=a+b$ and
$s=\operatorname{RN}(a+b)$.
\end{proof}
